\newpage
\thispagestyle{empty}

\begin{center}
    {\fontsize{14}{17}\selectfont \textbf{ADVERTENCIA}}
\end{center}

\vspace{1cm}

\noindent
La Universidad Santo Tomás no se hace responsable de las opiniones y conceptos expresados en el trabajo de grado, solo velará por qué no se publique nada contrario al dogma ni a la moral católica y porque el trabajo no tenga ataques personales y únicamente se vea el anhelo de buscar la verdad científica.

\newpage
\begin{center}
	{\large\bfseries DEDICATORIA}
\end{center}

\vspace{2cm}

\begin{flushright}
	\textit{
		A mi madre, el pilar fundamental de mi existencia y el motor inagotable que me impulsa a despertar cada día con el firme deseo de alcanzar mis sueños. Cada paso que he dado en este largo camino universitario lleva impreso tu esfuerzo, tu amor infinito y tu fe inquebrantable en mí. Este logro es tan tuyo como mío, porque sin tus sabios consejos, tus sacrificios silenciosos y tu abrazo cálido en los momentos de mayor duda, nada de esto habría sido posible. Gracias por darme la vida y las alas para volar hasta aquí.
	}
	
	\vspace{1cm}
	
	\textit{
		A Angie Marcela Díaz, por ser mi luz, mi refugio y mi paz; eres como ese amanecer que llega otra vez para disipar mis dudas. Por brindarme tu amor y tu apoyo incondicional en absolutamente cada aspecto de mi vida. Gracias por tomar mi mano en los días de mayor estrés durante el desarrollo de este proyecto, por celebrar mis pequeñas victorias cotidianas, por escucharme con infinita paciencia y por creer en mí incluso cuando yo mismo dudaba. Tu presencia ha hecho que este arduo camino de la ingeniería sea mucho más hermoso y llevadero. Esta meta alcanzada también lleva tu nombre.
	}
	
	\vspace{1cm}
	
	\textit{
		A mi hermano Antonio Camacho y su esposo Edilberto Pérez, quienes fueron el impulso inicial que necesitaba para dar el primer y más importante paso hacia la universidad. Su confianza en mí llegó incluso antes que la mía propia, y sin ese apoyo generoso y decisivo, este sueño tal vez nunca habría encontrado su punto de partida. Gracias por animarme a creer que este camino también era posible para mí.
	}
\end{flushright}

\newpage
\thispagestyle{empty}

\begin{center}
	{\large\bfseries AGRADECIMIENTOS}
\end{center}

\vspace{1cm}

\noindent
El desarrollo de este proyecto no habría sido posible sin el apoyo incondicional de personas excepcionales que creyeron en este proceso.

Agradezco profundamente al ingeniero Jaime Vitola Oyaga, director de este proyecto, por su respaldo institucional, por confiar en esta propuesta tecnológica y por defender con firmeza este trabajo frente a la comunidad académica.

Al ingeniero John Monzaide Álvarez Cely, codirector de esta investigación, le expreso mi más sincera gratitud por su paciencia, su orientación detallada y por guiarme paso a paso en cada fase técnica y metodológica.

Un agradecimiento sumamente especial al ingeniero Jesús Renato De La Rosa Martínez, asesor del proyecto. Gracias por todo el conocimiento técnico compartido de manera desinteresada y por ser una guía invaluable a lo largo del desarrollo de este sistema. Por encima de todo, valoro enormemente su sincera amistad y su constante acompañamiento personal y profesional.

Al ingeniero Carlos Javier Mojica Casallas, docente de la facultad, quiero expresarle un profundo agradecimiento por su disposición constante, sus valiosos consejos y toda la ayuda brindada a lo largo de mi formación académica; sus enseñanzas fueron piezas clave en mi recorrido como estudiante.

Un reconocimiento muy especial y profundo a la docente Elizabeth Mayorga, quien con su dedicación me transmitió el verdadero amor por las matemáticas y fue la brújula principal que me encaminó hacia la ingeniería. Sin aquellas primeras experiencias de aprendizaje tan excepcionales e inspiradoras de su mano, recorrer este camino y alcanzar esta meta simplemente no habría sido posible. De igual forma, extiendo mi sincera gratitud a la profesora Dora Téllez por sus invaluables enseñanzas, y a los docentes Ericilia Pedraza y Eddy Alexander Galviz por su apoyo incondicional y calidad humana, tanto dentro como fuera de las aulas durante mi etapa escolar.


Finalmente, extiendo mi entera gratitud al señor Frank Muller y a todo el equipo de trabajo del restaurante Gastro Uno 58 Bar. Su gran generosidad al abrir las puertas de su establecimiento y permitir la activa participación de sus empleados durante la prueba piloto fue una pieza fundamental para materializar este proyecto y llevarlo a la realidad operativa. Gracias por su absoluta confianza.
\vspace{1cm}


\newpage
{\bfseries Conflicto de Intereses}


\vspace{1cm}

\noindent
Los autores declaran no tener ningún conflicto de intereses. Los patrocinadores no participaron en el diseño del estudio, la recopilación, el análisis ni la interpretación de los datos publicados, la redacción del documento.
